%% ============================================================
%%  §8  TypeScript Extension and Cross-Language Unification
%% ============================================================

\section{TypeScript Extension and Cross-Language Unification}
\label{sec:typescript}

The TypeScript extension demonstrates that the \RustM monad and \ProgramSpec
specification language are genuinely language-agnostic.
\texttt{TypeScript/ElabSpec.lean} contains \textbf{33 theorems} (0 sorry)
covering 14 TypeScript functions — including a \texttt{while} loop — and
6 cross-language unification theorems that formally assert both a Rust and a
TypeScript implementation satisfy the \emph{same} specification.
An additional \textbf{6 theorems} in \texttt{TypeScript/BugDetection.lean}
demonstrate how the kernel acts as a bug oracle: it catches two classic
programming errors and confirms the fixes, all by \rfl.

\subsection{The Verified TypeScript Fragment}
\label{sec:ts-fragment}

TypeScript's execution model differs substantially from Rust's: dynamic
typing, garbage collection, prototype-based objects, no ownership.
The fragment we verify is a well-defined subset:

\begin{center}
\small
\begin{tabular}{@{}ll@{}}
\toprule
\textbf{Feature}                  & \textbf{Status} \\
\midrule
\texttt{number} (modelled as $\mathbb{Z}$) & Supported \\
\texttt{boolean}                   & Supported \\
\texttt{string} (constants and concatenation) & Partial (in evaluator; no theorems) \\
Arithmetic, comparison, logical ops & Supported \\
Strict equality (\texttt{===}), strict inequality (\texttt{!==}) & Supported \\
\texttt{if}/\texttt{else}, ternary \texttt{?:} & Supported \\
\texttt{while} loops with mutable variable update & Supported \\
\texttt{let}/\texttt{const} bindings & Supported \\
\texttt{typeof}                    & Supported \\
\midrule
Closures, arrow functions          & \textbf{Not supported} \\
Objects (\texttt{\{\}}), prototype chain & \textbf{Not supported} \\
Arrays (indexing in evaluator; no theorems) & Partial \\
\texttt{async}/\texttt{await}, Promises & \textbf{Not supported} \\
Type coercions (\texttt{==}, \texttt{+} on mixed types) & \textbf{Not supported} \\
IEEE~754 floating-point (NaN, Infinity, \texttt{-0}) & \textbf{Not supported} \\
\bottomrule
\end{tabular}
\end{center}

\noindent
The \texttt{number} type is modelled as $\mathbb{Z}$ (Lean's \texttt{Int})
rather than IEEE~754 \texttt{double}. This is a deliberate simplification:
integer arithmetic is exact under this model, which is sound for programs
that only use integer-valued expressions. Programs depending on floating-point
behaviour (NaN propagation, rounding, \texttt{-0}) are outside scope.

\subsection{The TypeScript AST and Evaluator}

\begin{lstlisting}[caption={TypeScript AST (actual definitions, selected).},
                   label={lst:tsast}]
inductive TSExpr where
  | lit   (l : TSLit)                            -- numeric/boolean literal
  | var   (idx : Nat)                            -- positional; idx 0 = first param
  | binOp (op : String) (l r : TSExpr)
  | unOp  (op : String) (e : TSExpr)
  | ite   (cond thenE elseE : TSExpr)

inductive TSStmt where
  | skip
  | return_ (e : TSExpr)
  | const   (name : String) (ty : Option TSTy) (e : TSExpr) (k : TSStmt)
  | set_    (idx : Nat) (e : TSExpr) (k : TSStmt)  -- mutable update env[idx]
  | ite     (cond : TSExpr) (thenB elseB : TSStmt)
  | while_  (cond : TSExpr) (body : TSStmt)
  | seq     (s1 s2 : TSStmt)
\end{lstlisting}

\noindent
\textbf{Key design choices:}
\begin{itemize}[leftmargin=2em]
  \item \textbf{Positional environments.} Variables are de Bruijn indices
        (\texttt{var idx}) into a \texttt{List TSValue}. This makes
        \texttt{evalExpr env (.var n) = env[n]?} kernel-transparent without
        any string equality, enabling all branching proofs to go through by
        \texttt{decide\_eq\_true/false} without sorry.

  \item \textbf{\texttt{const} + \texttt{set\_}.} Immutable bindings push a
        value to the front of the environment; \texttt{set\_} updates an existing
        slot in place. This pair covers both functional-style local declarations
        and imperative loop-variable mutation.

  \item \textbf{Same monad.} \texttt{evalTSFun} returns \texttt{RustM Unit TSValue}
        — exactly the type returned by \texttt{evalFun} for Rust. The common
        return type is what makes cross-language unification statements type-check.
\end{itemize}

The evaluator follows the same fuel-parametric, \texttt{match}-on-\texttt{Except}
discipline as the LLBC interpreter:

\begin{lstlisting}[caption={TypeScript evaluator entry point.},label={lst:tseval}]
def evalTSFun (f : TSFunDef) (fuel : Nat) (args : List TSValue) :
    RustM Unit TSValue :=
  fun _ =>
    let initEnv : Env := args     -- positional; idx 0 = first arg
    match evalStmt fuel f.body initEnv with
    | Except.ok (.return_ v, _) => Except.ok (v, ())
    | Except.ok (.next,     _)  => Except.error (.explicit "fell off end")
    | Except.error e            => Except.error e
    | _                         => Except.error (.explicit "break outside loop")
\end{lstlisting}

\subsection{Verified TypeScript Functions}

The 14 functions proved correct are:
\texttt{return42}, \texttt{id}, \texttt{neg}, \texttt{add}, \texttt{sub},
\texttt{max}, \texttt{isZero}, \texttt{mul}, \texttt{min},
\texttt{abs}, \texttt{not\_gate}, \texttt{clamp}, and \texttt{sum\_to}.
The first 12 are pure functions; \texttt{sum\_to} uses a \texttt{while} loop.

\paragraph{Pure functions.}
For functions with no branches (e.g., \texttt{add}, \texttt{neg}), the proof
is by \rfl — the kernel reduces the closed term to its value:

\begin{lstlisting}[caption={T3: neg proved by rfl.},label={lst:tsneg}]
def tsNegFun : TSFunDef := {
  body := .return_ (.unOp "-" (.var 0))  -- x = args[0]
  ...
}
theorem elab_ts_neg (x : Int) :
    evalTSFun tsNegFun 10 [.num x] at () |= .pureOutput (· = .num (-x)) :=
  ⟨.num (-x), (), rfl, rfl⟩
\end{lstlisting}

\paragraph{Branching functions.}
For functions with conditionals (\texttt{max}, \texttt{abs}, \texttt{clamp}),
we supply a branch hypothesis and use \texttt{decide\_eq\_true/false}
in a two-pass simp:

\begin{lstlisting}[caption={T10: abs (x $\ge$ 0 branch).},label={lst:tsabs}]
theorem elab_ts_abs_nonneg (x : Int) (h : x ≥ 0) :
    evalTSFun tsAbsFun 10 [.num x] at () |= .pureOutput (· = .num x) := by
  refine ⟨.num x, (), ?_, rfl⟩
  simp only [evalTSFun, tsAbsFun]
  simp only [evalStmt, evalExpr, evalBinOp, evalLit,
             List.getElem?_cons_zero, decide_eq_true h]
\end{lstlisting}

\paragraph{The while loop: sum\_to.}
\texttt{sum\_to(n)} computes $0 + 1 + \cdots + (n-1)$.
In TypeScript, this is a \texttt{while} loop over two mutable variables
\texttt{s} (accumulator) and \texttt{i} (counter):

\begin{lstlisting}[language=TypeScript,caption={TypeScript sum\_to source.},label={lst:tssumto}]
function sum_to(n: number): number {
  let s = 0;
  let i = 0;
  while (i < n) { s += i; i += 1; }
  return s;
}
\end{lstlisting}

The \texttt{TSFunDef} uses \texttt{const} to push \texttt{s} and \texttt{i}
onto the positional environment (making \texttt{i = env[0]},
\texttt{s = env[1]}, \texttt{n = env[2]}), and \texttt{set\_} to update
them in-place each iteration:

\begin{lstlisting}[caption={T14: tsSumToFun definition.},label={lst:tssumtofun}]
def tsSumToFun : TSFunDef := {
  name := "sum_to"; params := [⟨"n", .number, false⟩]; retTy := .number
  body :=
    .const "s" none (.lit (.num 0))     -- env = [0, n]
    (.const "i" none (.lit (.num 0))    -- env = [0, 0, n]; i=0,s=1,n=2
    (.seq
      (.while_ (.binOp "<" (.var 0) (.var 2))   -- while i < n
        (.set_ 1 (.binOp "+" (.var 1) (.var 0)) -- s = s + i
        (.set_ 0 (.binOp "+" (.var 0) (.lit (.num 1))) -- i = i + 1
          .skip)))
      (.return_ (.var 1))))             -- return s
  async_ := false
}
\end{lstlisting}

The correctness theorems are ground instances proved by \rfl (the kernel
evaluates the loop concretely):

\begin{lstlisting}[caption={T14 ground instances.},label={lst:tssumtothm}]
theorem elab_ts_sumTo_five :
    evalTSFun tsSumToFun 100 [.num 5] at () |= .pureOutput (· = .num 10) :=
  ⟨.num 10, (), rfl, rfl⟩

theorem elab_ts_sumTo_ten :
    evalTSFun tsSumToFun 100 [.num 10] at () |= .pureOutput (· = .num 45) :=
  ⟨.num 45, (), rfl, rfl⟩
\end{lstlisting}

\subsection{Cross-Language Unification Theorems (U1--U6)}
\label{sec:unification}

Six unification theorems establish that Rust and TypeScript implementations
satisfy the \emph{same} \ProgramSpec predicate:

\begin{lstlisting}[caption={U1: max — same spec from both languages.},
                   label={lst:unified}]
theorem unified_max_either (a b : Int) :
    (evalFun [] maxFun 10 [.int a, .int b] at () |=
      .pureOutput (fun v => v = .int a ∨ v = .int b)) ∧
    (evalTSFun tsMaxFun 10 [.num a, .num b] at () |=
      .pureOutput (fun v => v = .num a ∨ v = .num b)) :=
  ⟨elab_max_either a b, elab_ts_max_either a b⟩
\end{lstlisting}

\begin{lstlisting}[caption={U5: sum\_to — same spec from Rust loop and TypeScript while loop.},
                   label={lst:unifiedsumto}]
theorem unified_sumTo_ten :
    (evalFun [] sumToFun 1000 [.int 10] at () |=
      .pureOutput (· = .int 45)) ∧
    (evalTSFun tsSumToFun 100 [.num 10] at () |=
      .pureOutput (· = .num 45)) :=
  ⟨elab_sumTo_ten, elab_ts_sumTo_ten⟩
\end{lstlisting}

\begin{lstlisting}[caption={U6: neg — symbolic unification.},label={lst:unifiedneg}]
theorem unified_neg (x : Int) :
    (evalFun [] negFun 10 [.int x] at () |= .pureOutput (· = .int (-x))) ∧
    (evalTSFun tsNegFun 10 [.num x] at () |= .pureOutput (· = .num (-x))) :=
  ⟨elab_neg x, elab_ts_neg x⟩
\end{lstlisting}

\noindent
The full set of 6 unification theorems (U1--U6) covers
\texttt{max} (either branch), \texttt{add} (no-crash), \texttt{min} (either),
\texttt{mul} (no-crash), \texttt{sum\_to} (ground, n=10), and \texttt{neg} (symbolic).

These theorems are the formal realisation of the paper's core claim:
a single \ProgramSpec instance describes implementations in both languages.
The proofs are simple conjunctions of the individual results; the substantive
work is in establishing compatible specifications on each side.
From a software-engineering perspective, if a function is later ported from
Rust to TypeScript (or vice versa), the unification theorem formally certifies
that the rewrite preserves the specification.

\subsection{Bug Detection and Repair}
\label{sec:ts-bug-detection}

A key practical benefit of the deep-embedding approach is that the Lean kernel
acts as a \emph{bug oracle}: if a function definition contains an error, a proof
attempt immediately reveals the wrong computed value through a failed \rfl reduction.
\texttt{TypeScript/BugDetection.lean} demonstrates two representative bugs.

\paragraph{Bug B1: off-by-one in \texttt{sum\_to}.}
The intended function computes $0 + 1 + \cdots + (n-1)$.
The buggy version uses \texttt{<=} instead of \texttt{<} in the loop guard:

\begin{lstlisting}[language=TypeScript,caption={Buggy TypeScript sum\_to.},label={lst:buggysum}]
// BUG: condition should be i < n
while (i <= n) { s += i; i += 1; }
\end{lstlisting}

\noindent
This causes the loop to run for $i = 0, 1, \ldots, n$ (inclusive), computing
$n(n{+}1)/2$ instead of $n(n{-}1)/2$.
Attempting to prove \texttt{sum\_to\_buggy(5) = 10} fails:
\texttt{⟨.num 10, (), rfl, rfl⟩} does not type-check because the kernel
reduces the term to \texttt{.num 15}.
Theorem B1 instead \emph{documents the bug} by proving the wrong value:

\begin{lstlisting}[caption={B1: kernel confirms the off-by-one bug.},label={lst:b1}]
theorem ts_sumTo_buggy_five :
    evalTSFun tsSumToBuggyFun 100 [.num 5] at () |= .pureOutput (· = .num 15) :=
  ⟨.num 15, (), rfl, rfl⟩   -- 15 = 0+1+2+3+4+5 (includes 5; should be 10)
\end{lstlisting}

\noindent
Fixing the guard to \texttt{<} recovers the correct \texttt{tsSumToFun}, for which
\texttt{elab\_ts\_sumTo\_five} (T13) proves the result is \texttt{.num 10}.

\paragraph{Bug B2: wrong zero guard in \texttt{safeDiv}.}
A safe-division function should return $a/b$ when $b \neq 0$ and $0$
otherwise.  The buggy version guards on $b > 0$, silently returning $0$
for negative divisors:

\begin{lstlisting}[language=TypeScript,caption={Buggy TypeScript safeDiv.},label={lst:buggydiv}]
// BUG: guard should be b !== 0
function safeDivBuggy(a: number, b: number): number {
  if (b > 0) return a / b;
  return 0;
}
\end{lstlisting}

\noindent
Theorem B3 exposes the bug: the kernel evaluates \texttt{safeDivBuggy(10,~-2)}
to \texttt{0}, not \texttt{-5}:

\begin{lstlisting}[caption={B3: kernel exposes the wrong-guard bug.},label={lst:b3}]
theorem ts_safeDiv_buggy_negative_divisor :
    evalTSFun tsSafeDivBuggyFun 10 [.num 10, .num (-2)] at ()
      |= .pureOutput (· = .num 0) :=
  ⟨.num 0, (), rfl, rfl⟩    -- returns 0 for b=-2 (bug!)
\end{lstlisting}

\noindent
Fixing the guard to \texttt{!==} yields \texttt{tsSafeDivFun}.
Theorems B4--B6 verify all three cases: negative divisor returns $-5$,
positive divisor returns $3$, and zero divisor safely returns $0$
without panicking:

\begin{lstlisting}[caption={B4--B6: fixed safeDiv verified for all divisor signs.},label={lst:b46}]
theorem ts_safeDiv_correct_neg :
    evalTSFun tsSafeDivFun 10 [.num 10, .num (-2)] at () |= .pureOutput (· = .num (-5)) :=
  ⟨.num (-5), (), rfl, rfl⟩

theorem ts_safeDiv_correct_zero :
    evalTSFun tsSafeDivFun 10 [.num 10, .num 0] at () |= .pureOutput (· = .num 0) :=
  ⟨.num 0, (), rfl, rfl⟩
\end{lstlisting}

\noindent
All six bug-detection theorems are proved by \rfl.
The proof methodology is the same as for correctness proofs: the only
difference is that for the buggy theorems, the stated output value is the
\emph{wrong} answer, serving as machine-certified documentation of the defect.

\subsection{Absence of Sorry}

All 39 TypeScript theorems (33 in \texttt{ElabSpec.lean} + 6 in \texttt{BugDetection.lean})
are fully kernel-checked (0 sorry).
The key enabler is the \textbf{positional environment}: because
\texttt{env[n]?} reduces by \texttt{List.getElem?} without any
\texttt{@[extern]} string-equality instance, all evaluation equations
reduce under \texttt{simp}. Branching theorems supply a hypothesis
\texttt{h} and discharge the \texttt{decide} call via
\texttt{decide\_eq\_true h} or \texttt{decide\_eq\_false (not\_le.mpr h)}.
